\section{Tendencias clave de la GenAI empresarial}

\subsection{Democratización de la IA}

\begin{importantbox}{Los desarrolladores sustituyen a los expertos en IA}
Antes, construir aplicaciones de IA requería científicos de datos y expertos en ML; ahora cualquier desarrollador —incluso un estudiante de secundaria— puede construir aplicaciones de IA usando la API de un LLM. Esto amplió el grupo de desarrolladores de IA en más de 100 veces, y es la causa fundamental del crecimiento explosivo de las aplicaciones de IA.
\end{importantbox}

Detrás de la democratización de la IA hay dos fuerzas: la primera es que cada vez más modelos base abren sus API; la segunda es que las herramientas de bajo código o sin código encapsulan el diseño de prompts, la limpieza de datos y los procesos de Ops en flujos de trabajo visuales. Juntas, estas dos fuerzas trasladaron las aplicaciones de IA de un puñado de data scientist a un gran número de ingenieros de software y equipos de negocio.

\subsection{Disminución de la demanda de datos}

\begin{itemize}
    \item Antes, las empresas necesitaban meses o incluso años para recopilar datos etiquetados
    \item Los modelos grandes ya contienen un enorme conocimiento de preentrenamiento, por lo que la cantidad de datos necesaria para el fine-tuning se redujo drásticamente
    \item Las empresas pueden poner en marcha un proyecto de IA en cuestión de días, no de años
\end{itemize}

Reducir la demanda de datos no significa que se pueda omitir la gobernanza de datos. Las empresas aún necesitan establecer estrategias de «catálogo de datos», «puntuación de calidad de datos» y «acceso diferenciado» para evitar que el modelo absorba contenido no conforme durante el entrenamiento. Muchas empresas optan por validar primero con datos públicos y después hacer el ajuste fino con un lote pequeño de datos de alta calidad.

\subsection{Iteración rápida de modelos}

\begin{itemize}
    \item Cada pocas semanas se publica un nuevo modelo, y las clasificaciones cambian constantemente
    \item Las empresas se enfocan en la \textbf{plataforma} y no en un modelo individual: necesitan una infraestructura que permita cambiar de modelo con flexibilidad
    \item Hace 18 meses solo había un modelo de punta (GPT-4); ahora hay varias opciones equiparables
\end{itemize}

La iteración frecuente también trae consigo problemas de gestión de versiones. Una de las características de Vertex AI es que permite definir varias versiones (intercambiables) dentro del mismo pipeline de Agent y registrar los resultados de experimentos A/B. Así, incluso si alguna versión presenta un regress, es posible revertirla rápidamente.

\subsection{Convergencia multimodal}

\begin{itemize}
    \item Los modelos unimodales están siendo sustituidos por modelos multimodales (Gemini admite de forma nativa texto, imagen, audio y video)
    \item La demanda empresarial de entradas multimodales (gráficos en documentos, análisis de contenido de video) crece día con día
\end{itemize}

La convergencia multimodal implica atender simultáneamente problemas de sincronización de datos, alineación y resolución. Gemini construye un embedding space unificado entre voz, imagen y texto, de modo que un mismo flujo de entrada puede ser analizado a la vez por varios tokenizers, lo que evita rupturas entre modalidades.

\subsection{Descenso continuo de costo y latencia}

\begin{itemize}
    \item El costo de las llamadas a la API ha bajado alrededor de 1.5 órdenes de magnitud
    \item La reducción de la latencia hace posibles más aplicaciones en tiempo real
    \item Los modelos MoE dispersos (que solo activan una parte de los parámetros) reducen aún más el costo de la inferencia
    \item Hallazgo interesante: ejecutar ciertos modelos de API resulta más económico que ejecutar modelos de código abierto equivalentes (porque se ahorran los costos de operación)
\end{itemize}

Tras la baja de costos, las empresas prefieren la ruta de «algoritmo + hardware heterogéneo»: activar GPU de baja latencia en las horas pico, habilitar TPU/CPU más económicas en las horas de menor demanda, e incluso hacer offload de una parte de las solicitudes a un datacenter regional. Este enfoque exige que la capa de orchestration pueda percibir las características reales de cada solicitud.

\subsection{Resumen de la sección}

Tendencias clave de la GenAI empresarial: democratización de la IA, disminución del umbral de datos, iteración rápida de modelos, convergencia multimodal y descenso continuo de costos.

